\documentclass[11pt]{article}
\usepackage{amsmath,amssymb,hyperref,geometry}
\geometry{margin=1in}
\title{When a Change in Dock Access Makes a Warehouse Agreement Unstable}
\author{}
\date{}
\begin{document}
\maketitle

\paragraph{Result.} Two reached states have the same occupancies, incumbent allocation and payoffs, and the same number of slots licensed to each operator. Only one permits a profitable blocking coalition. The difference is whether a staging crew can legally move before its partner's collection slot. For this family the earliest attainable completion is an exact summary for strict blocking of the protected allocation, not for replaying or recovering every permitted deviation. A contingent agreement using the earliest feasible pairing eliminates the block without appropriating an old reservation.

\paragraph{Plant, authority, timing and information.}
Three independently operated groups $A,B,C$ share a civil warehouse transfer berth. At tick $-1$ the berth authority admits an old job for $A$. It starts at zero, uses the sole berth on $[0,1)$, and pays $7$ utility units upon physical release at tick $1$. No agreement may revoke, move, or compensate away the old occupancy. Utility is expressed in arbitrary integer units and is not estimated from warehouse operations. All actors know this law and the public access register. At the decision cut at tick zero, after old work starts, the authority announces the current sets $S_B,S_C$ of licensed positive start slots (the register can include slots beyond this service horizon) and a deadline $H$ fixed before optional access is allocated. Retain the original job and hold berth capacity fixed across these histories. Each group can only nominate its own licensed slots. The horizon ends at $H$; there are no future reputation payoffs, external transfers or later coalition rounds.

$B$ owns one optional batch. For a paired service, $B$ must stage it at the berth at $s\in S_B$ and $C$ collect at $t\in S_C$. Both actions occupy one tick; staging moves the batch into an adjacent non-berth buffer at $s+1$. A settling/handling rule requires $t-s\ge \ell$ for integer $\ell\ge 1$; the batch need not occupy the berth between operations. Collection completes at $e=t+1\le H$. The berth physically excludes overlapping starts. Release precedes the next start at the same tick. A bilateral, conditional acceptance at the decision cut binds both later actions for that pair and the payout if both finish. Once accepted, no further withdrawal is admissible until completion. Until acceptance either owner may withhold consent. The gate has no solo package mode: either both licensed nominations are accepted atomically or the sole new batch stays with $B$. This is an admission rule for an optional batch, not a license to stop old work.

\paragraph{Utilities and deviations.}
The incumbent offers only old service; its utility is $(7,0,0)$, so $0$ is each optional owner's fallback. At the pre-acceptance cut use strict blocking: a coalition blocks a proposed allocation if it can implement a legal alternative that strictly improves every member, holding the outsiders to their existing commitments. Singletons may abstain or nominate one leg; a nominee cannot compel an outsider to consent to a new package. A joint proposal by $B,C$ may select $(s,t)$ and may transfer integer credits $q$ from $C$ to $B$ with $|q|\le K$, charged only for actual paired completion and enforceable under the stated agreement. No transfers involving $A$ are permitted. Payoffs of the pair are
\[
 u_B(e,q)=\beta_B-\lambda_B e+q,\qquad
 u_C(e,q)=\beta_C-\lambda_C e-q,
\]
where each $\lambda_i$ is a nonnegative integer. The reward coefficients are net of the fixed effort for a completed pair; no part of a reward is earned by a sole unmatched nomination. Group $A$ receives $7$ in every admissible outcome. Thus no coalition containing $A$ can strictly block, and no singleton can strictly improve from the incumbent. These statements use the selected all-members-strict criterion, not weak improvements with a strict gain for only one member, and do not prove robustness to arbitrary side deals.

\paragraph{Matched failure and exact summary.}
With $S_B=\{1\}$, $S_C=\{2\}$, $\ell=1$ and $H=4$, the alternative starts staging at 1 after old release and collection at 2, finishing at 3. Swap the licenses to $S_B=\{2\}$ and $S_C=\{1\}$. The incumbent can run unchanged in both histories, producing the same physical custody and agreed utility. The roster, free berth intervals, slot counts, and union of licensed slots agree. Collection at 1 cannot follow staging at 2. The forward case blocks for $\beta_B=6,\beta_C=5,\lambda_B=\lambda_C=1,K=2$; e.g. $q=-2$ gives optional utilities $(1,4)$. No optional pair can run in the reverse case. In general, assigning two singleton licenses separated by at least $\ell$ before the deadline and swapping the assignment yields the same failure for every set of payoff parameters admitting strict gains at the forward completion.

Write $E=\infty$ if no valid pair exists. Otherwise $E=1+\min\{t\in S_C:t\ge\min S_B+\ell,\ t+1\le H\}$. The earliest staging choice leaves every collection option available that a later staging choice could have left, which proves the formula without assuming a joint-deviation oracle. Set $g_i(e)=\beta_i-\lambda_i e$. For finite $E$, exact blocking of the incumbent reduces to
\[
  \max(-K,1-g_B(E))\ \le\ \min(K,g_C(E)-1).
\]
Indeed these are the necessary and sufficient integer bounds for $u_B,u_C>0$ at the earliest completion. A later profitable pair would use a transfer $q$ that also makes both parties strictly better off at $E$, since neither untransferred reward increases with time. Conversely a valid $q$ at $E$ lifts to the actual licenses $s=\min S_B,t=E-1$, with old work still respected. For fixed lag and horizon but variable reward coefficients the quotient by $E$ is coarsest for these incumbent-blocking tests: if $e<e'$ take $\beta_B=\beta_C=e+1$ and $\lambda_B=\lambda_C=1$. Then the earlier pair gives each one credit at $q=0$, while the later pair has nonpositive total payoff (and $\infty$ admits no pair). This quotient does not preserve access order, intermediate deadlines, or the whole deviation relation. The lifting witness is checked against the actual access register, not recoverable from $E$ alone. With time-dependent transfer bounds, a reward for delay, occupied staging buffer, or a per-leg identity constraint, the proof must be revisited.

\paragraph{Repaired agreement.}
Publish the public access register and permit a contingent bilateral match after it is known, without modifying anyone's existing licenses or older commitments. If the criterion passes, offer the witnessed earliest pair with the least permitted $q$ in the feasible interval; otherwise retain the old allocation. In the positive case every member of the optional pair voluntarily accepts over its fallback, and a singleton cannot replace this matched joint proposal without the other party's assent. A different joint proposal cannot strictly improve both: total payoff at its completion cannot exceed the selected earliest total, and transfers conserve the sum. $A$'s earlier work is untouched. In the negative case the exact criterion rules out any feasible strict block. Thus this clearing rule selects a stable allocation at every modeled cut, including access changes that leave previously promised work intact. It neither preserves a guarantee to run an already accepted optional package across an arbitrary later revocation nor creates a physical access right from a payout.

\paragraph{Direct check and literature boundary.}
The matched direct approach enumerates all licensed starts and credit values without compressing the access history, then tests both utilities. Its proof is the same admissible-deviation law; the earliest-time summary is useful only under the named monotonicity and fixed-cap assumptions. No measured runtime, incentive compatibility under private valuations, empirical stability, or proof-maintenance advantage is asserted. The optional proposal changes the coalition structure at one public cut. \href{https://www.sciencedirect.com/science/article/pii/S016726810200015X}{Arnold and Schwalbe, \emph{Dynamic coalition formation and the core}, 2002} model ongoing coalition and payoff choices with adaptation and relate absorbing states to core allocations. \href{https://pages.nyu.edu/debraj/Papers/KonishiRay.pdf}{Konishi and Ray, \emph{Coalition Formation as a Dynamic Process}, 2003, \S1.1} let possible coalitional moves depend on state and compare continuation values. The state-dependent move correspondence here is in that spirit, but derives from exclusive physical use and time-bounded consent, not an abstract characteristic function. The one-cut finite certificate is not a dynamic equilibrium or convergence theorem, nor a priority claim.

\paragraph{Evidence.}
The direct scan compares 3,072 access/lag/deadline combinations and 331,776 utility cases (22,272 with a strict block), enumerating schedules and credits separately from the profile formula. The event interpreter additionally runs all 960 legal two-leg executions across those physical combinations and the four featured forward/reverse incumbent/repair traces. These checks supplement the general proof above. They do not establish deployment, acceptance beyond the declared finite envelope, or stability under other deviation rules.

\end{document}
